\documentclass[10pt]{article}
\usepackage{graphicx, libertine, multicol, setspace, siunitx, xcolor}
\usepackage[margin=1in]{geometry}
\definecolor{red1}{RGB}{160,0,20}
\usepackage[colorlinks = true, urlcolor = red1]{hyperref}
\usepackage[super]{nth}

\begin{document}
\noindent
\begin{center}
  \Large{\textbf{EC 201: Intermediate Microeconomic Analysis}} \\
  \medskip
  \large{\textbf{Boston University Summer Term II 2017}}
  \bigskip
\end{center}
\textbf{Class times:} Monday, Tuesday, Wednesday 6:00PM--8:30PM (July \nth{3} to August \nth{9}). \\
\textbf{Location:} \href{https://www.bu.edu/classrooms/classroom/cas-b20/}{CAS B20}, 725 Commonwealth Ave \\
\textbf{Instructor:} Christoph Walsh\\
\textbf{Office Hours:} Monday \& Wednesday 4:30PM--5:30PM \\
\textbf{Office Hours Location:}
\href{https://www.google.com/maps/place/Boston+University+School+of+Social+Work/@42.350845,-71.1052384,15z/data=!4m2!3m1!1s0x0:0xe9b0a401410ded28?sa=X&ved=0ahUKEwibjuTl_b7MAhXKJR4KHdfDAScQ_BIIdDAO}{SSW} B16, 264 Bay State Road \\
\textbf{Email:} \href{mailto:walshcb@bu.edu}{walshcb@bu.edu} \\
\textbf{Course Website:}
\href{https://walshc.github.io/ec201}{walshc.github.io/ec201}
\subsection*{\textsc{Course Summary}}
The goal of the course is to examine the topics studied in Introductory Microeconomics (EC101) in more detail and with more rigor. This will continue to develop your economic intuition and expand your knowledge of the analytical tools useful in modeling consumer and firm behavior. A solid foundation in this course will greatly assist you in future economics courses, as the tools used here are building blocks for the tools you will use in the future.

We will begin by studying the foundations of the demand curve through individual consumption choice, along with intertemporal and risky choice.  Analogously, we'll consider firms' choices of inputs and production quantities.  Together, these will allow us to consider market equilibria. We will then study strategic interaction and use it to analyze various other market structures.  Lastly we will look at settings where the market fails due to asymmetric information or failures of rationality.

\subsection*{\textsc{Textbook and Materials}}
The main notes for the course will be written on the whiteboard during class time. Problem sets, practice exams (and their solutions) and extra notes will be posted on the course website. The textbook for the course is:

\medskip
\textsc{Varian, Hal.~R.} (2014) \emph{Intermediate Microeconomics with Calculus: A Modern Approach}, \nth{1} edition, W.~W.~Norton \& Company
\medskip

\noindent While the textbook is not required, the material covered in the course and the material in the book are very similar and reading it may be a useful addition to your class notes.

\subsection*{\textsc{Prerequisites}}
A sound understanding of the topics covered in EC101 will be crucial as we will be discussing some of those topics in greater detail. We will also use basic calculus, algebra, and graphing techniques very frequently, which I will review on the first day. If you are uncomfortable with any of the topics covered on the first day, it is very important that you come see me as soon as possible.

\subsection*{\textsc{Grading}}
\subsubsection*{Attendance (10\%)}
Attendance is extremely important. There are very few lectures relative to a regular term course and we will be covering a lot of material in each lecture. Thus, missing a class could result in falling behind rather quickly.

\subsubsection*{Class Participation (10\%)}
Long class times in the evenings make it difficult to stay focused. Interaction and class participation can make this easier. A portion of this grade will go to students who participate well during the course. I will also ask students to solve homework problems on the board during class time.

\subsubsection*{Midterm Exam (35\%)}
The midterm exam will be held on Monday July \nth{24} during class time. It will include all content covered up to that date.

\subsubsection*{Final Exam (45\%)}
The final exam will be held on Wednesday August \nth{9} during class time. It will be \textbf{cumulative}, including all content covered in the class since the first day. There will be no make-up exams for either exam except in extreme cases.

\subsection*{\textsc{Problem Sets}}
The midterm and final exams will consist of problems that you will need to solve and the problem sets will serve as practice for the exams. Doing problems to practice the techniques is the only way to master the material in preparation for exams. The problem sets will also serve as a way to emphasize the important concepts that we discuss in class, as well as give you an opportunity to explore some issues at a deeper level. There will be one problem set each week and we will spend a portion of the class on Monday going over the previous week's problem set.

\subsection*{\textsc{Other Issues}}
I will be teaching the course on the whiteboard because I believe that for this kind of material (involving lots of graphs and mathematics) this will be the easiest for students to follow. Therefore please do not use laptops to take notes in class.  It is very difficult and also distracting to transfer the mathematics and graphs to the computer, so your notes will suffer. I also don't recommend using a tablet with a stylus as the graphs often require a ruler and are quite detailed. It is also good practice to draw the graphs on paper as that is how you will do it in the exams.

If you have a disability that allows you extra time on exams or any other accommodations, please let me know on the first day and, as soon as possible, give me a note from BU disability services. If any students need extra time I will need to book the classroom for a longer time period in advance which takes time.

The BU academic conduct code is in effect during this course. You can
find the details spelled out here if you are not familiar:
\href{http://www.bu.edu/academics/policies/academic-conduct-code/}
{bu.edu/academics/policies/academic-conduct-code/}

\subsection*{\textsc{Course Outline}}
The following lists the topics we will cover and their associated textbook chapters:
\begin{multicols}{2}
  \begin{itemize}
    \item Mathematics Review (Mathematical Appendix)
    \item The Budget Constraint (Ch.~2)
    \item Preferences and Utility (Ch.~3-4)
    \item Choice and Demand (Ch.~5-6)
    \item Intertemporal Choice (Ch.~10)
    \item Uncertainty (Ch.~12)
    \item Technology and Profit Maximization (Ch.~19-20)
    \item Cost Curves (Ch.~22)
    \item Firm Supply and Industry Supply (Ch.~23-24)
    \item Equilibrium (Ch.~16)
    \item Demand Elasticities (Ch.~15)
    \item Monopoly (Ch.~25-26)
    \item Game Theory (Ch.~29-30)
    \item Oligopoly (Ch.~28)
  \end{itemize}
\end{multicols}
\end{document}
